\subsection{三维卷积的存储读取比}
\label{l11:subsec:three_dimensional_reads}

考虑步幅为 $1$ 的三维卷积，核的三个维度为 $5\times5\times7$，权重存于常量内存。一个线程块负责 $16\times16\times16$ 的逻辑输出块，并先将完整的联合输入区域装入共享内存，其中每个输入元素只从全局内存读取一次。卷积阶段，对每个输出的每个核位置执行一次共享内存标量读取。要求内部输出块的全局读取次数与共享读取次数之比，忽略权重访问和输出写入。

设输出块的三个维度为 $T_x,T_y,T_z$，卷积核的对应维度为 $K_x,K_y,K_z$。沿每个轴，相邻输出窗口的联合范围都比输出范围多出 $K_i-1$ 个输入位置。因此，装入共享内存的联合输入区域包含
\begin{equation}
  G=(T_x+K_x-1)(T_y+K_y-1)(T_z+K_z-1)
  \label{l11:eq:three_dim_global}
\end{equation}
个元素。每个元素只装入一次，故这也是全局读取次数。一个输出需要读取 $K_xK_yK_z$ 个共享输入，输出总数为 $T_xT_yT_z$，所以共享读取次数为
\begin{equation}
  S=T_xT_yT_zK_xK_yK_z.
  \label{l11:eq:three_dim_shared}
\end{equation}

代入本题参数，输入区域尺寸为 $(16+4)\times(16+4)\times(16+6)=20\times20\times22$，故 $G=8\,800$。逻辑输出数为 $16^3=4\,096$，每个输出读取 $5\times5\times7=175$ 个共享输入，故 $S=4\,096\times175=716\,800$。所求比值为
\begin{equation}
  G:S=8\,800:716\,800=11:896
  \approx \boxed{1:81.45}.
  \label{l11:eq:three_dim_ratio}
\end{equation}
这里比较的是规定计算方式下的标量读取次数。共享内存读取具有不同于全局内存读取的成本，实际实现还包含同步、地址计算和算术操作，因此该比值不是运行时间的加速比。

逻辑输出块的尺寸也不等于线程块必须采用的尺寸。$16^3$ 个输出可以由较少线程分担，一个线程计算多个输出；实际线程块必须满足设备的线程数和资源限制。输出的分工方式应与上面采用的读取计数模型保持一致。

同样的方法适用于长方体输出块。例如，若第二个输出维度为 $161$，即输出块为 $16\times161\times16$，输入区域变为 $20\times165\times22$。此时
\begin{equation}
\begin{aligned}
  G&=72\,600,\\
  S&=16\cdot161\cdot16\cdot175=7\,212\,800,\\
  G:S&=72\,600:7\,212\,800\approx1:99.35.
\end{aligned}
\label{l11:eq:three_dim_rectangular}
\end{equation}
输出块的具体尺寸决定了晕区开销，不能在读取比计算中省略任一轴的长度。

\subsection{多通道卷积的浮点运算总数}
\label{l11:subsec:flop_problem}

一幅 $100\times200$ 的彩色图像含 $3$ 个输入通道。卷积层使用 $9\times9$ 核生成 $10$ 个输出特征图，每个输出特征图都结合全部 $3$ 个输入通道，步幅为 $1$，不进行填充。一次乘加计为 $2$ 次浮点运算，仅统计卷积，忽略偏置、激活、池化和其他操作。

首先计算有效输出范围。高度方向的窗口左上角可从 $0$ 移动到 $100-9=91$，因此共有 $92$ 个位置；宽度方向可从 $0$ 移动到 $200-9=191$，因此共有 $192$ 个位置。每个输出特征图有 $92\times192=17\,664$ 个值，全部十个特征图共有 $176\,640$ 个输出值。

每个输出值分别与三个通道中的 $9\times9$ 窗口做乘积并累加，需要 $3\times9\times9=243$ 次乘加。总乘加数为
\begin{equation}
  N_{\mathrm{MAC}}
  =10\cdot92\cdot192\cdot3\cdot9^2
  =42\,923\,520.
  \label{l11:eq:flop_problem_mac}
\end{equation}
按一次乘加计 $2$ FLOPs 的约定，总浮点运算数为
\begin{equation}
  N_{\mathrm{FLOP}}
  =2\cdot42\,923\,520
  =\boxed{85\,847\,040\ \text{FLOPs}}.
  \label{l11:eq:flop_problem_result}
\end{equation}
若错用输入的 $100\times200$ 个空间位置，就会把不属于有效卷积输出的边界位置也计入；若漏掉输入通道因子 $3$，则只计算了一个输入通道的贡献。输出通道数 $M$ 决定需要多少组核，输入通道数 $C$ 决定每个输出内积的长度，这两个因子承担不同作用。

该层只有 $10\times3\times9^2=2\,430$ 个卷积权重，却执行了四千多万次乘加，因为每个权重都在 $92\times192$ 个输出位置复用。权重共享降低参数量，同时保留了对所有输出位置逐一完成计算的要求。

\subsection{计算组织的统一关系}
\label{l11:subsec:synthesis}

卷积的数学依赖可以分为两组索引：$(b,m,h,w)$ 确定独立输出，$(c,p,q)$ 确定该输出内部的乘加归约。直接实现保留这组张量索引；矩阵实现把 $(c,p,q)$ 合并为 $r$，把 $(h,w)$ 合并为 $\ell$。两种表示的核心量对应为
\begin{equation}
  D=CK^2,\qquad L=H_{\mathrm{o}}W_{\mathrm{o}},\qquad
  N_{\mathrm{output}}=BML,\qquad
  N_{\mathrm{MAC}}=BMDL.
  \label{l11:eq:core_relations}
\end{equation}
共享内存复用相邻输出的重叠输入，矩阵分块沿输出通道复用输入，寄存器子块在单个线程内复用操作数。显式与隐式展开保持上述索引和输出含义；融合卷积、池化、偏置与激活时，还必须保持各阶段的运算次序、布局和同步条件。
